\documentclass{article}
\usepackage{amsmath,amssymb,amsthm}
\usepackage{algorithm}
\usepackage[noend]{algpseudocode}
\usepackage{fullpage}
\newtheorem{theorem}{Theorem}
\newtheorem{lemma}{Lemma}
\newtheorem{assumption}{Assumption}
\newcommand{\gap}{\mathcal{G}}
\newcommand{\diam}{\mathrm{diam}}
\begin{document}

\section*{Algorithm Name}
\textbf{Frank–Wolfe (Conditional Gradient) Method for Non‑Convex Smooth Optimization}

\section*{Mathematical Formulation}
Let $f:\mathbb{R}^d\rightarrow\mathbb{R}$ be differentiable and $C\subseteq\mathbb{R}^d$ a non‑empty compact convex set.
Given a step‑size sequence $\{\gamma_k\}_{k\ge 0}\subset(0,1]$, the iterates $\{x_k\}$ are generated by  

\begin{align}
s_k &:= \arg\min_{s\in C}\;\langle \nabla f(x_k),\,s\rangle , \label{eq:LMO}\\
x_{k+1} &:= (1-\gamma_k)x_k + \gamma_k s_k . \label{eq:update}
\end{align}

The quantity 
\[
\gap(x):=\max_{s\in C}\langle -\nabla f(x),\,s-x\rangle
      =\langle \nabla f(x),\,x-s^\star(x)\rangle ,
\qquad s^\star(x):=\arg\min_{s\in C}\langle \nabla f(x),s\rangle,
\]
is called the \emph{Frank–Wolfe (FW) gap} and serves as a stationarity measure.

\section*{Pseudocode}
\begin{algorithm}[H]
\caption{Frank–Wolfe for Non‑Convex $f$}
\begin{algorithmic}[1]
\Require Differentiable $f$, compact convex $C$, initial $x_0\in C$, stepsizes $\{\gamma_k\}$
\For{$k=0,1,2,\dots$}
    \State Compute gradient $g_k \gets \nabla f(x_k)$
    \State \textbf{Linear minimization oracle:}
          \[s_k \gets \arg\min_{s\in C}\langle g_k,s\rangle\] \Comment{Eq.~\eqref{eq:LMO}}
    \State Update
          \[x_{k+1} \gets (1-\gamma_k)x_k + \gamma_k s_k\] \Comment{Eq.~\eqref{eq:update}}
    \If{termination criterion met (e.g., $\gap(x_k)\le\epsilon$)}
        \State \textbf{break}
    \EndIf
\EndFor
\end{algorithmic}
\end{algorithm}

\section*{Dry Run Example}
Consider the univariate quadratic 
\[
f(w) = (w-3)^2 ,\qquad C=[0,5],\qquad w_0=0,\qquad \gamma_k\equiv \eta=0.1 .
\]

\begin{tabular}{c|c|c|c}
Iteration $k$ & $g_k=\nabla f(w_k)$ & $s_k=\arg\min_{s\in[0,5]}\langle g_k,s\rangle$ & $w_{k+1}=(1-\eta)w_k+\eta s_k$ \\ \hline
0 & $g_0 = 2(w_0-3)= -6$ & $s_0 = 5$ (since $g_0<0$) & $w_1 = 0.9\cdot0+0.1\cdot5=0.5$ \\ 
  & $f(w_0)=9$ &  & $f(w_1)=(0.5-3)^2=6.25$ \\ \hline
1 & $g_1 = 2(0.5-3)= -5$ & $s_1 = 5$ & $w_2 = 0.9\cdot0.5+0.1\cdot5 =0.95$ \\ 
  & $f(w_1)=6.25$ &  & $f(w_2)=(0.95-3)^2\approx4.2025$ \\ \hline
2 & $g_2 = 2(0.95-3)= -4.1$ & $s_2 = 5$ & $w_3 = 0.9\cdot0.95+0.1\cdot5 \approx1.355$ \\ 
  & $f(w_2)\approx4.2025$ &  & $f(w_3)=(1.355-3)^2\approx2.701$ 
\end{tabular}

The FW gap at each iterate equals $\gap(w_k)=\langle g_k, w_k-s_k\rangle$, e.g. $\gap(w_0)=(-6)(0-5)=30$, decreasing as the algorithm proceeds.

\section*{Assumptions}
\begin{assumption}[Smoothness]\label{asm:smooth}
$f$ is $L$‑smooth on an open set containing $C$, i.e.
\[
\|\nabla f(x)-\nabla f(y)\|\le L\|x-y\|,\qquad \forall x,y\in C .
\]
\end{assumption}
\begin{assumption}[Bounded Domain]\label{asm:bounded}
$C$ is compact with finite diameter
\[
D:=\diam(C)=\max_{x,y\in C}\|x-y\|<\infty .
\]
\end{assumption}
\begin{assumption}[Access to Exact LMO]\label{asm:lmo}
At each iteration the linear minimization oracle returns an exact solution $s_k$ of \eqref{eq:LMO}.
\end{assumption}
\begin{assumption}[Step‑size]\label{asm:stepsize}
The step‑sizes satisfy $\gamma_k\in(0,1]$ and
\[
\sum_{k=0}^{\infty}\gamma_k = \infty,\qquad
\sum_{k=0}^{\infty}\gamma_k^2 < \infty .
\]
A convenient choice is $\gamma_k = \frac{2}{k+2}$.
\end{assumption}

\section*{Main Convergence Theorem}
\begin{theorem}\label{thm:main}
Under Assumptions \ref{asm:smooth}–\ref{asm:stepsize}, the iterates $\{x_k\}$ generated by \eqref{eq:update} satisfy
\[
\min_{0\le k\le T-1}\mathbb{E}\big[\gap(x_k)\big]
   \;\le\; \frac{2L D^2}{\sum_{k=0}^{T-1}\gamma_k}
   \;=\; O\!\left(\frac{1}{\sqrt{T}}\right) ,
\]
when the step‑size sequence $\gamma_k = \frac{c}{\sqrt{k+1}}$ with a constant $c>0$ is used.
Consequently,
\[
\lim_{T\to\infty}\min_{0\le k<T}\gap(x_k)=0 .
\]
\end{theorem}

\section*{Theoretical Properties}
\begin{itemize}
\item \textbf{Stability.} The update \eqref{eq:update} is a convex combination of two points in $C$, thus $x_{k+1}\in C$ for all $k$, guaranteeing feasibility without any projection.
\item \textbf{Hyper‑parameter sensitivity.} The convergence rate depends only on the product $L D^2$ and the decay of $\gamma_k$. Faster decays (e.g. $\gamma_k=2/(k+2)$) yield an $O(1/T)$ bound on the averaged gap, while the $1/\sqrt{T}$ schedule is robust to noise in the gradient.
\item \textbf{Comparison to SGD.} Unlike SGD, the method never leaves $C$, and the stationarity measure is the FW gap rather than gradient norm. The $O(1/\sqrt{T})$ rate matches the optimal rate for first‑order methods on non‑convex smooth problems without variance reduction.
\end{itemize}

\section*{Discussion}
The theorem shows that even without convexity the Frank–Wolfe algorithm drives the FW gap to zero, implying that every limit point is a \emph{stationary point} of the constrained problem $\min_{x\in C}f(x)$. The dependence on the domain diameter $D$ is intrinsic to conditional‑gradient methods because the linear oracle can only move at most a distance $D$ per iteration.

\section*{Preliminaries (Definitions and Lemmas)}
\begin{lemma}[Descent Lemma for $L$‑smooth functions]\label{lem:descent}
For any $x,y\in C$,
\[
f(y) \le f(x) + \langle \nabla f(x),y-x\rangle + \frac{L}{2}\|y-x\|^2 .
\]
\end{lemma}
\begin{proof}
Standard; follows from integrating the Lipschitz condition on the gradient. \qedhere
\end{proof}

\begin{lemma}[Bound on distance moved]\label{lem:stepbound}
The update \eqref{eq:update} satisfies
\[
\|x_{k+1}-x_k\| = \gamma_k\|s_k-x_k\|\le \gamma_k D .
\]
\end{lemma}
\begin{proof}
Since $s_k,x_k\in C$, $\|s_k-x_k\|\le D$. Multiplying by $\gamma_k$ yields the claim. \qedhere
\end{proof}

\section*{Main Proof (Step‑by‑Step Derivation)}
We prove Theorem \ref{thm:main}. All non‑trivial steps are numbered.

\bigskip
\noindent\textbf{Step 1.} Apply Lemma \ref{lem:descent} with $x=x_k$ and $y=x_{k+1}$:
\[
f(x_{k+1}) \le f(x_k) + \langle \nabla f(x_k),x_{k+1}-x_k\rangle + \frac{L}{2}\|x_{k+1}-x_k\|^2 .
\tag{1}
\]

\noindent\textbf{Step 2.} Substitute the update rule $x_{k+1}-x_k = \gamma_k (s_k-x_k)$:
\[
\langle \nabla f(x_k),x_{k+1}-x_k\rangle = \gamma_k\langle \nabla f(x_k),s_k-x_k\rangle .
\tag{2}
\]

\noindent\textbf{Step 3.} Use the definition of the FW gap:
\[
\gap(x_k)=\max_{s\in C}\langle -\nabla f(x_k),s-x_k\rangle
        = -\langle \nabla f(x_k),s_k-x_k\rangle .
\tag{3}
\]
Hence $\langle \nabla f(x_k),s_k-x_k\rangle = -\gap(x_k)$.

\noindent\textbf{Step 4.} Combine (1)–(3) and Lemma \ref{lem:stepbound}:
\[
\begin{aligned}
f(x_{k+1}) &\le f(x_k) -\gamma_k \gap(x_k) 
            + \frac{L}{2}\gamma_k^2\|s_k-x_k\|^2 \\[4pt]
           &\le f(x_k) -\gamma_k \gap(x_k) + \frac{L D^2}{2}\gamma_k^2 .
\end{aligned}
\tag{4}
\]

\noindent\textbf{Step 5.} Rearrange (4) to isolate the gap term:
\[
\gamma_k\gap(x_k) \le f(x_k)-f(x_{k+1}) + \frac{L D^2}{2}\gamma_k^2 .
\tag{5}
\]

\noindent\textbf{Step 6.} Sum (5) over $k=0,\dots,T-1$:
\[
\sum_{k=0}^{T-1}\gamma_k\gap(x_k) 
   \le f(x_0)-f(x_T) + \frac{L D^2}{2}\sum_{k=0}^{T-1}\gamma_k^2 .
\tag{6}
\]

\noindent\textbf{Step 7.} Since $f$ is lower bounded on the compact set $C$, let $f_{\inf}:=\min_{x\in C}f(x)$.
Then $f(x_T)\ge f_{\inf}$, yielding
\[
\sum_{k=0}^{T-1}\gamma_k\gap(x_k) 
   \le f(x_0)-f_{\inf} + \frac{L D^2}{2}\sum_{k=0}^{T-1}\gamma_k^2 .
\tag{7}
\]

\noindent\textbf{Step 8.} Divide both sides of (7) by $\sum_{k=0}^{T-1}\gamma_k$:
\[
\frac{\sum_{k=0}^{T-1}\gamma_k\gap(x_k)}{\sum_{k=0}^{T-1}\gamma_k}
   \le \frac{f(x_0)-f_{\inf}}{\sum_{k=0}^{T-1}\gamma_k}
      + \frac{L D^2}{2}\frac{\sum_{k=0}^{T-1}\gamma_k^2}{\sum_{k=0}^{T-1}\gamma_k} .
\tag{8}
\]

\noindent\textbf{Step 9.} The left‑hand side is a convex combination of the gaps, therefore
\[
\min_{0\le k\le T-1}\gap(x_k) 
   \le \frac{\sum_{k=0}^{T-1}\gamma_k\gap(x_k)}{\sum_{k=0}^{T-1}\gamma_k} .
\tag{9}
\]

\noindent\textbf{Step 10.} Choose the step‑size $\gamma_k = \frac{c}{\sqrt{k+1}}$ with $c>0$.
Then
\[
\sum_{k=0}^{T-1}\gamma_k 
   \ge c\int_{0}^{T}\frac{1}{\sqrt{t}}dt = 2c\sqrt{T},
\qquad
\sum_{k=0}^{T-1}\gamma_k^2 
   = c^2\sum_{k=0}^{T-1}\frac{1}{k+1}
   \le c^2(1+\log T) .
\tag{10}
\]

\noindent\textbf{Step 11.} Plug (10) into (8)–(9). Using the bound $\frac{\sum\gamma_k^2}{\sum\gamma_k}\le \frac{c(1+\log T)}{2\sqrt{T}}=O(T^{-1/2})$ we obtain
\[
\min_{0\le k<T}\gap(x_k)
   \le \frac{f(x_0)-f_{\inf}}{2c\sqrt{T}}
      + \frac{L D^2}{2}\cdot O\!\left(\frac{1}{\sqrt{T}}\right)
   = O\!\left(\frac{1}{\sqrt{T}}\right) .
\tag{11}
\]

\noindent\textbf{Step 12.} Since the right‑hand side tends to zero as $T\to\infty$, the minimum FW gap over the first $T$ iterates converges to zero, completing the proof. \qed

\section*{Rate Derivation (With Constants)}
From (11) we can write an explicit constant:
\[
\min_{0\le k<T}\gap(x_k)
   \le \frac{f(x_0)-f_{\inf}}{2c\sqrt{T}}
      + \frac{L D^2 c(1+\log T)}{4\sqrt{T}} .
\]
Choosing $c = \sqrt{f(x_0)-f_{\inf}}/(L D^2)^{1/2}$ balances the two terms and yields
\[
\min_{k<T}\gap(x_k) \le \frac{2\sqrt{L D^2\big(f(x_0)-f_{\inf}\big)}}{\sqrt{T}} .
\]

\section*{Special Cases}
\begin{itemize}
\item \textbf{Convex $f$.} The same derivation gives the classic $O(1/T)$ rate for the primal gap $f(x_k)-f^\star$ when the step‑size $\gamma_k=2/(k+2)$ is used.
\item \textbf{Exact line‑search.} If $\gamma_k$ is chosen by exact line‑search, i.e. 
\[
\gamma_k = \arg\min_{\gamma\in[0,1]} f\big((1-\gamma)x_k+\gamma s_k\big),
\]
the inequality in Step 4 becomes an equality, and the bound improves to $O(1/T)$ even without convexity.
\item \textbf{Stochastic gradients.} When only an unbiased estimator $\tilde g_k$ of $\nabla f(x_k)$ is available, the same proof holds in expectation provided $\mathbb{E}\|\tilde g_k-\nabla f(x_k)\|^2\le\sigma^2$ and the LMO uses $\tilde g_k$.
\end{itemize}

\end{document}